\scanpage{140}\mspage{71}

donc l'opération de $(\mathbb Z/n\mathbb Z)^*$ sur $\overrightarrow{\mathrm{Pol}}(E)$ passe au quotient en une opération sur
\[
\deleted{\mathrm{Pol}(E)}\qquad\mathrm{Pol}(E)=\overrightarrow{\mathrm{Pol}}(E)/(\omega\mapsto\omega^{\unsure{\nu|-1}}),
\]
enfin l'élément $-1\in(\mathbb Z/n\mathbb Z)^*$ opère trivialement sur $\mathrm{Pol}(E)$, en vertu de 35 a), d'où l'opération (32).

Pour $\nu\in(\mathbb Z/n\mathbb Z)^*/\pm1=Z'_n$, et $\alpha$ structure $n$-polygonale sur un ens $E$, on désigne par $\alpha^\nu$ la structure \deleted{\unsure{induite}}\inserted{transformée} par $\nu$. Ceci fixé, j'en viens à l'axiome massue :

% slanted
\margin{Si $A\subset\mathfrak P_2(E)$ définit la structure polygonale, on désigne par $A^{(\nu)}\subset\mathfrak P_2(E)$ sa transformée par $\nu\in Z'_n$.}

(33)\qquad (To 6)\qquad $\forall\lambda\in Z$, \deleted{$\unsure{s\in S},\ \unsure{\varepsilon}$} \inserted{$\exists\nu=\nu(\lambda)\in\underbrace{(\mathbb Z/n\mathbb Z)^*/\pm1}_{Z'_n}$} (évidemment unique) tel que pour tout $s\in S$, la structure $n$-polygonale \deleted{induite} $\alpha_{\lambda.i_0,\lambda.\unsure{s|b}}$ induite par $A_{\lambda.i_0}=A_\lambda$\inserted{ (plutôt : $T_\lambda i_0$)} sur $\mathrm{Et}_{A_0}(s)=\mathrm{Et}_{A_\lambda}(\lambda s)$, soit égale à $(\alpha_{i_0,s})^\nu$.

Cela détermine donc la structure $A_\lambda$ en termes de $\nu=\nu(\lambda)\in Z'_n$, \inserted{notons-la $A_0^{(\nu)}$,} donc dans le cas où $P_0$ n'est pas l'octaèdre (\deleted{\unsure{au cas}}\inserted{dans lequel} de toutes façons $Z=\{1,\unsure{\underline{\circ}|\underline{-}|\varepsilon}\}$ opère trivialement sur $\Pi$, réduit à un pt !), cela détermine aussi $\lambda$ en termes de $\nu$ :

% archive sheet number: 139
